\section{Sisteme de Recomandare în Domeniul Modei}
Sistemele de recomandare (RS) reprezintă coloana vertebrală a personalizării în aplicațiile moderne. În industria modei, acestea trebuie să depășească barierele algoritmilor clasici, adaptându-se la natura extrem de vizuală și subiectivă a preferințelor vestimentare. Această secțiune analizează paradigmele existente și justifică alegerea filtrării bazate pe conținut pentru proiectul StyleGenAI.

\subsection{Categorii de Sisteme de Recomandare}
Conform taxonomiei stabilite de Adomavicius și Tuzhilin \cite{adomavicius2005toward}, sistemele de recomandare se împart în trei categorii majore, fiecare având aplicabilitate distinctă în modă:

\begin{itemize}
    \item \textbf{Filtrarea Colaborativă (Collaborative Filtering)}: Se bazează pe "înțelepciunea mulțimii". Dacă utilizatorul A și utilizatorul B au gusturi similare în trecut, sistemul recomandă lui A ce a plăcut lui B.
    \item \textbf{Filtrarea Bazată pe Conținut (Content-Based)}: Recomandă articole similare cu cele pe care utilizatorul le deține deja, analizând atributele produsului (culoare, textură, formă).
    \item \textbf{Sisteme Hibride}: Combină cele două abordări pentru a elimina dezavantajele fiecăreia.
\end{itemize}

\subsection{De ce Filtrarea Colaborativă eșuează în Garderoba Personală}
Deși filtrarea colaborativă domină platformele de streaming (Netflix, Spotify), ea este ineficientă pentru aplicația noastră din cauza problemei \textit{Cold Start}.
Într-un dulap virtual, utilizatorul își încarcă propriile haine (poze unice). Aceste articole nu există în dulapurile altor utilizatori, deci nu există un istoric comun de interacțiuni ("Like-uri" de la alți oameni) pe baza căruia să se facă recomandări. 
Astfel, soluția optimă pentru StyleGenAI este \textbf{Filtrarea Bazată pe Conținut Vizual}, care analizează haina însăși, nu cine a mai purtat-o.

\subsection{Algoritmul de Similaritate și Extragerea Trăsăturilor}
Pentru a recomanda o piesă vestimentară care se "potrivește" cu alta, sistemul trebuie să înțeleagă similaritatea vizuală. Aceasta nu se face pe pixelii bruți ai imaginii, ci pe \textit{vectori de trăsături} (embeddings) extrași printr-o Rețea Neuronală Convoluțională (CNN).

Formal, măsura similarității între două articole $A$ și $B$, reprezentate prin vectori de trăsături $\vec{a}$ și $\vec{b}$, se calculează folosind \textbf{Similaritatea Cosinus}. Aceasta măsoară cosinusul unghiului dintre cei doi vectori într-un spațiu multidimensional:

\begin{equation}
\text{similarity}(\vec{a}, \vec{b}) = \cos(\theta) = \frac{\vec{a} \cdot \vec{b}}{\|\vec{a}\| \|\vec{b}\|}
\end{equation}

Unde:
\begin{itemize}
    \item Dacă scorul este 1, vectorii sunt identici (haine foarte similare).
    \item Dacă scorul este 0, vectorii sunt ortogonali (haine complet diferite).
    \item Dacă scorul este -1, vectorii sunt opuși.
\end{itemize}

\subsection{Implementare Practică în Python}
În arhitectura backend a StyleGenAI, acest concept matematic este implementat folosind biblioteca \texttt{scikit-learn}. Codul de mai jos ilustrează funcția care compară amprentele vizuale a două articole vestimentare:

\begin{verbatim}
import numpy as np
from sklearn.metrics.pairwise import cosine_similarity

def calculate_style_match(embedding_top, embedding_bottom):
    """
    Calculează gradul de potrivire dintre un tricou și o pereche de pantaloni.
    Args:
        embedding_top (np.array): Vectorul trăsăturilor vizuale pentru partea de sus.
        embedding_bottom (np.array): Vectorul trăsăturilor pentru partea de jos.
    Returns:
        float: Scor de compatibilitate (0.0 - 1.0).
    """
    # Redimensionare necesară pentru funcția library-ului
    vec_a = embedding_top.reshape(1, -1)
    vec_b = embedding_bottom.reshape(1, -1)
    
    # Calculăm similaritatea și normalizăm rezultatul
    raw_score = cosine_similarity(vec_a, vec_b)[0][0]
    return max(0.0, raw_score) # Asigurăm că nu returnăm valori negative
\end{verbatim}

Această funcție este esențială pentru "generatorul de ținute", permițând sistemului să filtreze rapid mii de combinații posibile și să le prezinte utilizatorului doar pe cele coezive vizual.

\subsection{Abordări bazate pe Secvențe (RNN)}
Dincolo de simpla similaritate între două piese, o ținută completă este o secvență: \textit{Top + Bottom + Shoes}. Cercetările recente, precum cele ale lui Han et al. \cite{han2017learning}, propun utilizarea rețelelor recurente (LSTM - Long Short-Term Memory) pentru a modela această relație.
Deși StyleGenAI folosește primar abordarea bazată pe reguli de stil și similaritate vizuală datorită eficienței computaționale pe dispozitive mobile, arhitectura permite extinderea viitoare către un model secvențial complet.

